%% LyX 2.4.4 created this file.  For more info, see https://www.lyx.org/.
%% Do not edit unless you really know what you are doing.
\documentclass[english]{article}
\usepackage[T1]{fontenc}
\usepackage[latin9]{inputenc}
\usepackage{units}
\usepackage{enumitem}
\usepackage{amsmath}
\usepackage{amsthm}
\usepackage{cancel}
\usepackage{geometry}
\geometry{verbose,tmargin=1in,bmargin=1in,lmargin=1in,rmargin=1in}

\makeatletter
%%%%%%%%%%%%%%%%%%%%%%%%%%%%%% Textclass specific LaTeX commands.
\numberwithin{equation}{section}
\numberwithin{figure}{section}
\newlength{\lyxlabelwidth}      % auxiliary length 

%%%%%%%%%%%%%%%%%%%%%%%%%%%%%% User specified LaTeX commands.
\usepackage{fancyhdr}
\usepackage{lastpage}
\pagestyle{fancy}
\usepackage{enumitem}
\usepackage{ifthen}
\everymath=\expandafter{\the\everymath\displaystyle}
%\DeclareMathSizes{10}{10}{10}{10}
\usepackage{multicol}

\usepackage{tikz}
\usetikzlibrary{patterns}
\usetikzlibrary{plotmarks}
\usepackage{pgfplots}
\pgfplotsset{compat=newest}

\usepackage{advdate}    % Advancing/saving dates
\usepackage{datetime}   % Dates formatting
\usepackage{datenumber} % Counters for dates
\usdate
% Added by lyx2lyx
\setlength{\parskip}{\smallskipamount}
\setlength{\parindent}{0pt}

\makeatother

\usepackage{babel}
\begin{document}
\renewcommand{\author}{Dr.\,James Ashe}

\newcommand{\classNum}{232}

\newcommand{\classSec}{A and B}

\newcommand{\examType}{Exam 2}

\renewcommand{\date}{\formatdate{11}{3}{2013}}

%%First page's header%%%%%%%%%%%%%%%%%%%%%%
\fancypagestyle{firststyle} %%header settings for first pagr
{
	\fancyhead[L]{MTH \classNum{}(\classSec{}) \author{}\\
	\textbf{\examType{}} \date{}}
	\fancyhead[C]{\textbf{Name:}}
	\setlength{\headsep}{14pt}
}
%%%%%%%%%%%%%%%%%%%%%%%%%%%%%%%%%%%%%%%%%%

%%Next page's header%%%%%%%%%%%%%%%%%%%%%%%
\setlist{nolistsep}
\lhead{\classNum{}(\classSec{})} 
\chead{\textbf{}} 
\cfoot{\thepage{}~of~\pageref{LastPage}} 
\setlength{\headsep}{5pt} 
%%%%%%%%%%%%%%%%%%%%%%%%%%%%%%%%%%%%%%%%%%%

%%Various settings; see the preamble for other settings%%%%%
%\setenumerate[0]{leftmargin=12pt,itemindent=0pt,labelwidth=0pt}
\setlist{topsep=.5pt}
\renewcommand{\labelenumi}{\textbf{\arabic{enumi}.}}
\renewcommand{\labelenumii}{\textbf{\alph{enumii}.}}
\thispagestyle{firststyle}
%%%%%%%%%%%%%%%%%%%%%%%%%%%%%%%%%%%%%%%%%%

\newcommand{\xmax}{0}
\newcommand{\xmin}{-\xmax}
\newcommand{\xstep}{0}
\newcommand{\ymax}{0}
\newcommand{\ymin}{-\ymax}
\newcommand{\ystep}{0} 

\textbf{Justify your answers and show your work. }Unless stated otherwise,
all problems are worth 5 points. 

\emph{Evaluate the following definite integral(s).}
\begin{enumerate}[resume]
\item ${\displaystyle \int_{-1}^{2}x^{2}(x^{3}+1)^{4}\,dx}$

Let $u=x^{3}+1$ so $dx=\frac{du}{3x^{2}}$, $u(-1)=0$, and $u(2)=9$.
\begin{eqnarray*}
{\displaystyle } & = & \int_{0}^{9}\cancel{x^{2}}u^{4}\frac{du}{3\cancel{x^{2}}}\\
 & = & \frac{1}{15}u^{5}|_{0}^{9}=\frac{9^{5}}{15}-0=\frac{9^{5}}{15}\\
 & = & \frac{19683}{5}=3936.6
\end{eqnarray*}
\vfill
\item ${\displaystyle \int_{-2}^{1}\frac{t}{\sqrt{4-t^{2}}}\,dt}$

Let $u=4-t^{2}$ so $dt=-\frac{du}{2t}$, $u(-2)=0$, and $u(1)=3$.
\end{enumerate}
\begin{eqnarray*}
 & = & {\displaystyle \int_{0}^{3}\cancel{t}\cdot u^{-\nicefrac{1}{2}}\frac{du}{-2\cancel{t}}}\\
 & = & -\frac{1}{2}\left(\frac{2}{1}\right)u^{\nicefrac{1}{2}}|_{0}^{3}=-\sqrt{3}\,\,(\approx-1.73205)
\end{eqnarray*}
\vfill

\emph{A full 200 gasoline tank begins to leaks out of at the rate
of $50-2t$, where t is measured in minutes.}
\begin{enumerate}[resume]
\item How much gasoline leaked out in the first $4$ minutes?

\[
{\displaystyle \int_{0}^{4}50-2t\,dt=\left(50t-\frac{2}{2}t^{2}\right)|_{0}^{4}=200-16=184\mbox{ gallons}}
\]
\vfill
\end{enumerate}
\emph{The function $v(t)=-\sin(t)$, on $0\leq t\leq\nicefrac{3}{2}\pi$,
is the velocity in m/sec of a particle moving along the x-axis.}
\begin{enumerate}[resume]
\item {[}2 points{]} If the particle has an initial position of $s(0)=1$,
find the function $s(t)$ giving its position.

\begin{eqnarray*}
s(t) & = & {\displaystyle -\int\sin(t)\,dt}\\
 & = & \cos(t)+C,\mbox{ }s(0)=1\Rightarrow C=0\\
 & = & \cos(t)
\end{eqnarray*}
\vfill
\item Find the total distance traveled over the given interval. 

\begin{multicols}{2}

$-\sin(t)=0$ on $[0,\nicefrac{3}{2}\pi]$ at $t=0$ and $t=\pi$.
So, 
\begin{eqnarray*}
\mbox{distance} & = & {\displaystyle \int_{0}^{\nicefrac{3}{2}\pi}\left|-\sin t\right|\,dt}\\
 & = & {\displaystyle \int_{0}^{\pi}\sin t\,dt+{\displaystyle \int_{\pi}^{\nicefrac{3}{2}\pi}-\sin t}\,dt}\\
 & = & \left(\left(1\right)-\left(-1\right)\right)+\left(0-\left(-1\right)\right)=3
\end{eqnarray*}
\columnbreak

\renewcommand{\xmax}{5}
\renewcommand{\xmin}{-.25}
\renewcommand{\xstep}{0} %must be xmin+n
\renewcommand{\ymax}{1}
\renewcommand{\ymin}{-1}
\renewcommand{\ystep}{0} %must be ymin+n
%%%%%%%
\begin{tikzpicture}
\begin{axis}[height=4cm, 
	width=5cm, 
	axis lines=middle,
	%axis line style={-latex}, 
	xmin=\xmin,xmax=\xmax, 
	ymin=\ymin,ymax=\ymax,
	%restrict y to domain=-1:10,
	%no markers, 
	xlabel={$x$},ylabel={$y$}, 
	%every axis x label/.append style={anchor=north}, 
	%every axis y label/.append style={anchor=east}, 
	area style,tension=0.08]
	
	\addplot+[domain=0:4.71239,fill,black!30!white]{-1*sin(deg(x))} \closedcycle;
	
	\addplot[domain=\xmin:\xmax,thick,blue]{-1*sin(deg(x))} node[left] {$$}; 
	
	%\addplot[dashed,thick,blue] coordinates{(.5,0) (.5,4)};
	
	%\node at (axis cs:.5,2) {$R$};
\end{axis}
\end{tikzpicture}\end{multicols}\newpage
\end{enumerate}
\emph{Use calculus to find the area of the region described.}
\begin{enumerate}[resume]
\item The region bounded by $y=x$ and $y=x^{2}-4x$.

\begin{eqnarray*}
R & = & {\displaystyle \int_{0}^{5}x-\left(x^{2}-4x\right)\,dx}\\
 & = & \frac{5}{2}x^{2}-\frac{1}{3}x^{3}|_{0}^{5}=\frac{125}{6}\approx20.8333
\end{eqnarray*}
\renewcommand{\xmax}{6}
\renewcommand{\xmin}{-6}
\renewcommand{\xstep}{0} %must be xmin+n
\renewcommand{\ymax}{8}
\renewcommand{\ymin}{-8}
\renewcommand{\ystep}{0} %must be ymin+n
%%%%%%%
\begin{tikzpicture}
\begin{axis}[height=4cm, 
	width=4cm, 
	axis lines=middle,
	%axis line style={-latex}, 
	xmin=\xmin,xmax=\xmax, 
	ymin=\ymin,ymax=\ymax,
	%restrict y to domain=-1:10,
	%no markers, 
	xlabel={$x$},ylabel={$y$}, 
	%every axis x label/.append style={anchor=north}, 
	%every axis y label/.append style={anchor=east}, 
	area style,tension=0.08]
	
	\addplot+[domain=0:5,fill,black!30!white]{x)} \closedcycle;
	\addplot+[domain=0:5,fill,black!30!white]{x^2-4*x} \closedcycle;
	\addplot+[domain=4:5,fill,white]{x^2-4*x} \closedcycle; 
	
	\addplot[domain=\xmin:\xmax,thick,smooth,red,samples=500]{x^2-4*x} node[left] {$x^2-4x$}; 
	\addplot[domain=\xmin:\xmax,blue,thick] {x} node[left] {$x$}; 	
	
	\node at (axis cs:2,0) {$R$};
\end{axis}
\end{tikzpicture}\vfill
\end{enumerate}
\emph{Find the definite integrals(s) which give the area of the region
described. }\textbf{\emph{For this problem you only need to write
the simplified integral which gives the area.}}
\begin{enumerate}[resume]
\item The region in the first quadrant bounded by $-3x+6$, $4x-4$, and
$-6x+6$. Use calculus

$4x-4=-3x+6\Rightarrow x=\frac{10}{7}$
\begin{eqnarray*}
R & = & {\displaystyle \int_{0}^{1}-3x+6-\left(-6x+6\right)\,dx+{\displaystyle \int_{1}^{\nicefrac{10}{7}}-3x+6-\left(4x-4\right)\,dx}}\\
 & = & {\displaystyle \int_{0}^{1}3x\,dx+{\displaystyle \int_{1}^{\nicefrac{3}{2}}-7x+10\,dx}}
\end{eqnarray*}
\renewcommand{\xmax}{2}
\renewcommand{\xmin}{-.5}
\renewcommand{\xstep}{0} %must be xmin+n
\renewcommand{\ymax}{8}
\renewcommand{\ymin}{-1.5}
\renewcommand{\ystep}{0} %must be ymin+n
%%%%%%%
\begin{tikzpicture}
\begin{axis}[height=4cm, 
	width=5cm, 
	axis lines=middle,
	%axis line style={-latex}, 
	xmin=\xmin,xmax=\xmax, 
	ymin=\ymin,ymax=\ymax,
	%restrict y to domain=-1:10,
	%no markers, 
	xlabel={$x$},ylabel={$y$}, 
	%every axis x label/.append style={anchor=north}, 
	%every axis y label/.append style={anchor=east}, 
	area style,tension=0.08]
	
	\addplot+[domain=0:1.5,fill,black!30!white]{-3*x+6} \closedcycle;
	\addplot+[domain=0:1,fill,white]{-6*x+6}  \closedcycle;
	\addplot+[domain=1:1.5,fill,white]{4*x-4} \closedcycle; 
	
	\addplot[domain=\xmin:\xmax,thick,black]{-3*x+6} node[left] {$$}; 
	\addplot[domain=\xmin:\xmax,thick,red]{4*x-4} node[left] {$$}; 
	\addplot[domain=\xmin:\xmax,blue,thick]{-6*x+6} node[left] {$$}; 	
	
	\addplot[dashed,thick,blue] coordinates{(1,0) (1,3)};
	
	\node at (axis cs:1,1.25) {$R$};
\end{axis}
\end{tikzpicture}\vfill
\end{enumerate}
\emph{Use the general slicing method and to find the volume of the
solid. }
\begin{enumerate}[resume]
\item {[}8 points{]} The solid lies between planes perpendicular to the
\emph{x}-axis at $x=0$ and $x=\nicefrac{\pi}{4}$. The cross-sections
are squares which are perpendicular to the \emph{x}-axis and have
a height of $\sqrt{\cos2x}$ measured positively from the \emph{x}-axis.

The cross-sections are square with sides of $\sqrt{\cos2x}$. So then
the area of each cross-section is $\left(\sqrt{\cos2x}\right)^{2}=\cos2x$.
\begin{eqnarray*}
V & = & {\displaystyle \int_{0}^{\nicefrac{\pi}{4}}\cos2x\,dx}\\
 & = & \frac{1}{2}\sin2x|_{0}^{\nicefrac{\pi}{4}}=\frac{1}{2}(1)-\frac{1}{2}0=\frac{1}{2}\mbox{ cubic units}
\end{eqnarray*}
\vfill\newpage
\end{enumerate}
\emph{Let $R$ be the region bounded by the following curves. Use
the disk or washer method to find the volume of the solid generated
when $R$ is revolved about the $x$-axis. }
\begin{enumerate}[resume]
\item $y=\sqrt{x}$, $x=0$, and $x=9$.

\begin{eqnarray*}
V & = & {\displaystyle \pi\int_{0}^{9}\left(\sqrt{x}\right)^{2}\,dx}\\
 & = & \pi\frac{1}{2}x^{2}|_{0}^{9}=\frac{\pi81}{2}\approx127.235\mbox{ cubic units}
\end{eqnarray*}
\vfill
\end{enumerate}
\textbf{\emph{For this problem You only need to write the simplified
integral which gives the volume.}}
\begin{enumerate}[resume]
\item $y=-16x$ and $y=8x^{2}+x^{3}$.

$-16x=8x^{2}+x^{3}\Rightarrow x\left(-x^{2}-8x-16\right)=x\left(x+4\right)\left(x+4\right)=0$.
So $x=0,-4$

\[
\pi{\displaystyle {\displaystyle \int}_{-4}^{0}\left[\left(-16x\right)^{2}-\left(8x^{2}+x^{3}\right)^{2}\right]\,dx=\frac{32768}{35}}\approx3.1019
\]
\vfill
\end{enumerate}
\emph{Find the length of the curve.}
\begin{enumerate}[resume]
\item $y=2x^{\nicefrac{3}{2}}$ from $x=0$ to $x=\frac{5}{4}$.

$y'=3x^{\nicefrac{1}{2}}$ and $y'{}^{2}=9x$. 

${\displaystyle \int_{0}^{\nicefrac{5}{4}}\sqrt{1+9x}}\,dx$ let $u=1+9x$.
So then $dx=\frac{du}{9}$.
\begin{eqnarray*}
L & = & {\displaystyle \frac{1}{9}\int_{0}^{\nicefrac{5}{4}}u^{\nicefrac{1}{2}}}\,du\\
 & = & \frac{1}{9}\frac{2}{3}u^{\nicefrac{3}{2}}|_{0}^{\nicefrac{5}{4}}=\frac{2}{27}\left(1+9\left(\frac{5}{4}\right)\right)^{\nicefrac{3}{2}}=\frac{335}{108}
\end{eqnarray*}
\vfill
\end{enumerate}
\emph{BONUS: Use calculus and the general slicing method to find a
general formula for the volume of a pyramid with a square base.}

Place the pyramid on the Cartesian plane such that the base is perpendicular
to the \emph{y-}axis and the peak touches the \emph{x}-axis (as we
did on the quiz). Let $s$ be length of the base, and $h$ be the
height. The equation which gives the distance from the \emph{x}-axis
to the side of the pyramid is the line $y=-\frac{sx}{2h}+\frac{b}{2}$
(the line has slope $\frac{\nicefrac{s}{2}-0}{0-h}$ and \emph{y}-intercept
$(0,\nicefrac{b}{2})$). So the area of the cross-section is given
by 
\[
A(x)=\left(2\left(-\frac{sx}{2h}+\frac{s}{2}\right)\right)^{2}=\left(\frac{s(h-x)}{h}\right)^{2}=\frac{s^{2}}{h^{2}}\left(h-x\right)^{2}
\]
and 
\begin{eqnarray*}
{\displaystyle V=\int_{0}^{h}\frac{s^{2}}{h^{2}}(h-x)^{2}\,dx} & = & \frac{s^{2}}{h^{2}}\int_{0}^{h}\left(h-x\right)^{2}\,dx\\
 & = & s^{2}\frac{1}{3h^{2}}\left(h-x\right)^{3}|_{0}^{h}=s^{2}\frac{h^{3}}{3h^{2}}=\frac{s^{2}h}{3}
\end{eqnarray*}

The area of the base, \textbf{$b=s^{2}$}. So this is often written
$V=\frac{1}{3}bh$.
\end{document}
